\section{Conclusiones}

\begin{enumerate}[leftmargin=*]

\item El modelo FOPDT identificado con el ensayo escalón,
$\Gpnom(s) = 0{,}979\,e^{-8{,}00 s}/(179{,}1 s + 1)$, representa adecuadamente la planta en la
región de operación del controlador, con un ajuste del \SI{88.7}{\percent} sobre datos no empleados
en su identificación, frente al \SI{74.4}{\percent} del modelo multinivel. Su limitación es
igualmente clara: el residuo de validación no es aleatorio sino creciente con el nivel de potencia,
de \SI{+0.88}{\celsius} a \SI{+2.20}{\celsius} entre $Q_1 = \SI{20}{\percent}$ y
\SI{50}{\percent}, lo que indica que una sola ganancia no describe todo el rango de la planta.

\item El anti-windup, y no la sintonización, fue lo que determinó el cumplimiento de las
especificaciones. Con el mismo $\Kp$ y el mismo $\Ti$ obtenidos por SIMC, la configuración sin
protección del integrador produce $M_p = \SI{30.1}{\percent}$ e incumple el requisito
$M_p \leq \SI{5}{\percent}$ por un factor de seis, mientras que con \emph{clamping} se obtiene
$M_p = \SI{0}{\percent}$ y $t_s = \SI{205}{\second}$. Haber aplicado correctamente las reglas SIMC
no demuestra por sí mismo que el controlador cumpla los requisitos.

\item Los dos diseños PI cumplen las especificaciones en MIL, pero con compromisos opuestos y sin
que ninguno domine al otro: el PI--SIMC establece \SI{47}{\second} antes, y el PI--CSD demanda un
comando inicial de \SI{108}{\percent} frente a \SI{274}{\percent} y mantiene el actuador saturado
\SI{21}{\second} frente a \SI{66}{\second}, además de contar con márgenes de estabilidad mayores.
La selección depende del criterio que se privilegie y no se deduce de las métricas de desempeño por
sí solas.

\item Las verificaciones SIL y PIL resultaron en igualdad exacta de $Q_1$ y $T_1$ sobre
\num{1801} muestras, lo que acredita que ni la generación de código ni su ejecución sobre el
Arduino Uno R4 Minima alteraron el comportamiento del controlador validado en MIL. Ese resultado no
constituye una validación del sistema completo: durante ambas pruebas la planta, el retardo y la
temperatura ambiente permanecen simulados y no se ejercitan los canales físicos de la TCLab, de
modo que el valor de estas etapas está en acotar las causas posibles de una discrepancia
posterior, no en anticiparla.

\item La implementación sobre la TCLab cumple los requisitos ---$t_s = \SI{298}{\second}$ y
$M_p = \SI{1.96}{\percent}$--- con un error estacionario de \SI{0.04}{\celsius}, muy inferior al
escalón mínimo del convertidor (\SI{0.489}{\celsius}). Las diferencias respecto de MIL se explican
principalmente por un error de ganancia del modelo: la planta requiere \SI{40.3}{\percent} de
potencia en régimen frente al \SI{36.8}{\percent} previsto, lo que corresponde a una ganancia
estática aparente cerca de un \SI{15}{\percent} menor que la nominal y explica el tiempo de
establecimiento un \SI{45}{\percent} mayor.

\end{enumerate}
